\answer{ex:11:1}{(а)~$T=k^2/\bigl(\Phi(\Delta I/I)^2\bigr)=9\,\text{с}$. (б)~$45$ датчиков, пятно $\approx7\times7$, разрешение хуже в $\approx7$ раз.}
\answer{ex:11:2}{(а)~$2$--$3.3$ бита на импульс, $22$--$34\,\%$ границы ($9.1$--$9.8$ бита). (б)~$\log_2\binom{400}{20}\approx111$ бит; в среднем $1$ общий тип.}
\answer{ex:11:3}{(а)~$\sigma/9\le I\le9\sigma$, $1.9$ порядка. (б)~$6$ и $7$.}
\answer{ex:11:4}{(а)~$f<343/0.44\approx780$~Гц. (б)~$625$ импульсов, $125$ нейронов.}
\answer{ex:11:5}{$21$ бит на событие, $4.8\cdot10^5$ событий/с, $7$ событий на пиксель края: $136$ пикселей/с. Обычная камера --- $1.25$ кадра/с.}
\answer{ex:11:6}{В логарифмах $\approx1.0\,\text{с}$ в обе стороны (теория $0.96\tau$); линейно $0.2\,\text{с}$ вверх и $3.0\,\text{с}$ вниз (теория $0.18\tau$ и $3.0\tau$).}
\answer{ex:11:7}{(а)~$\ln I$ распределён логистически с медианой $\ln\sigma$ и масштабом $1$: разброс $\pi/\sqrt3$ в натуральных логарифмах, $0.79$ порядка. (б)~$r=r_{\max}\Phi_I(I)^2$.}
